\chapter{Pendahuluan}
\section{Latar Belakang}
\par
Keterikatan kuantum (\textit{quantum entanglement}) merupakan salah satu fenomena fundamental dalam mekanika kuantum yang berperan penting dalam perkembangan teori dan aplikasi informasi kuantum, seperti komputasi kuantum, komunikasi kuantum, dan kriptografi kuantum \cite{horodecki2009quantum}. Dalam konteks sistem multipartit, keterikatan multipartit sejati (\textit{genuine multipartite entanglement}) menjadi perhatian utama karena mencerminkan korelasi nonklasik yang melibatkan seluruh subsistem secara simultan dan tidak dapat direduksi menjadi kombinasi keterikatan bipartit \cite{guhne2009entanglement}. Oleh karena itu, pemahaman dan karakterisasi keterikatan multipartit sejati memiliki signifikansi teoretis dan aplikatif yang tinggi\cite{santra2024genuine,sarkar2018multipartite}.
\par
Seiring dengan kemajuan penelitian eksperimental yang berhasil merealisasikan keterikatan multipartit sejati, tantangan teoretis yang berkaitan dengan pendefinisian dan pengukuran derajat keterikatan tersebut masih menjadi permasalahan yang terbuka \cite{guhne2009entanglement}. Berbeda dengan sistem bipartit yang telah memiliki ukuran keterikatan yang relatif mapan, pengukuran keterikatan multipartit sejati menghadapi kompleksitas matematis yang lebih tinggi\cite{physreva_measurement_2015,nature_entanglement_2024}. Berbagai ukuran keterikatan yang telah diusulkan dalam literatur menunjukkan keterbatasan tertentu, baik dari sisi sensitivitas maupun keterterapannya, khususnya ketika diterapkan pada sistem dengan banyak subsistem atau pada keadaan campuran \cite{eltschka2014quantifying}.
\par
Pendekatan geometris terhadap pengukuran keterikatan kuantum merupakan salah satu alternatif yang berkembang untuk mengatasi permasalahan tersebut \cite{wei2003geometric}. Pendekatan ini memanfaatkan struktur geometris ruang keadaan kuantum untuk mengkarakterisasi keterikatan, sehingga memungkinkan deskripsi yang lebih sistematis terhadap hubungan antara besaran fisik dan struktur matematis keadaan kuantum. Dalam kerangka ini, pengisian konkurensi (\textit{concurrence fill}) diperkenalkan sebagai ukuran geometris yang dirancang untuk mengukur keterikatan multipartit sejati pada keadaan murni, dengan mengaitkan derajat keterikatan pada luas area segitiga konkurensi yang mendasarinya \cite{2021triangle}.
\par
Pada sistem tiga qubit, pengisian konkurensi dapat diformulasikan dalam kaitannya dengan besaran keterikatan lain yang telah dikenal, seperti \textit{three-tangle} dan \textit{partial tangles} \cite{coffman2000distributed}. Formulasi ini menghasilkan representasi matematis yang lebih ringkas dan mudah dioperasionalkan, sehingga memungkinkan analisis yang lebih terstruktur terhadap berbagai kelas keadaan multipartit sejati. Melalui pendekatan tersebut, karakteristik keterikatan pada keadaan kelas GHZ (Greenberger--Horne--Zeilinger) dan kelas W dapat dianalisis secara sistematis, yang merupakan dua kelas fundamental keterikatan multipartit sejati dalam sistem tiga qubit \cite{dur2000three}.
\par
Selain keadaan murni, perhatian juga diberikan pada keadaan campuran, mengingat sistem kuantum dalam kondisi realistis umumnya berada dalam keadaan campuran akibat interaksi dengan lingkungan. Untuk kasus tertentu, seperti campuran peringkat-2 dari keadaan kelas GHZ dan W yang digeneralisasi, struktur keterikatan menunjukkan sifat yang menarik, di mana tiga konkurensi dapat bernilai nol pada wilayah tertentu dalam ruang keadaan yang membentuk suatu polihedron nol (\textit{zero polyhedron}) \cite{aggarwal2025classification}. Analisis terhadap kondisi ini memungkinkan penentuan pengisian konkurensi untuk eigenstate dari campuran yang bersesuaian serta penetapan batas atas pengisian konkurensi pada keadaan campuran dengan konkurensi nol.
\par
Berdasarkan uraian tersebut, studi ini difokuskan pada analisis teoretis pengisian konkurensi sebagai ukuran keterikatan multipartit sejati pada sistem tiga qubit. Penelitian ini bertujuan untuk memperdalam pemahaman terhadap formulasi matematis pengisian konkurensi, mengkaji karakteristiknya dalam mengklasifikasikan kelas keadaan multipartit sejati, serta menelaah penerapannya pada keadaan murni dan campuran. Diharapkan hasil studi ini dapat memberikan kontribusi dalam memperjelas peran pengisian konkurensi dalam teori keterikatan multipartit serta memperkaya kajian teoretis di bidang fisika kuantum dan informasi kuantum.

\section{Rumusan Masalah}
Berdasarkan latar belakang yang telah diuraikan, maka rumusan masalah dalam penelitian ini adalah sebagai berikut:
\begin{enumerate}
    \item Bagaimana pengisian konkurensi dapat diformulasikan dan digunakan sebagai ukuran geometrik untuk mengidentifikasi keterbelitan multipartit sejati pada sistem tiga qubit?
    \item Bagaimana kriteria klasifikasi keadaan tiga qubit ke dalam kelas GHZ dan W dapat ditentukan berdasarkan karakteristik pengisian konkurensi?
\end{enumerate}

\section{Tujuan Penelitian}
Adapun tujuan dari penelitian ini adalah sebagai berikut:
\begin{enumerate}
    \item Mengkaji dan merumuskan penggunaan pengisian konkurensi sebagai ukuran geometrik dalam mengidentifikasi dan mengklasifikasikan keterbelitan kuantum pada sistem tiga qubit, khususnya keterbelitan multipartit sejati.
    \item Menentukan kriteria klasifikasi keadaan tiga qubit ke dalam kelas GHZ dan W berdasarkan karakteristik pengisian konkurensi.
\end{enumerate}

\section{Batasan Masalah}
Agar pembahasan dalam penelitian ini lebih terarah dan terfokus, maka batasan masalah yang digunakan adalah sebagai berikut:
\begin{enumerate}
    \item Sistem kuantum yang dikaji dibatasi pada sistem tiga qubit.
    \item Analisis keterikatan multipartit sejati difokuskan pada ukuran pengisian konkurensi.
    \item Keadaan kuantum yang dikaji meliputi keadaan murni serta keadaan campuran tertentu, khususnya campuran peringkat-2 dari keadaan kelas GHZ dan W.
    \item Pendekatan yang digunakan bersifat teoretis dan analitis, tanpa melibatkan simulasi numerik maupun verifikasi eksperimental.
\end{enumerate}

\section{Manfaat Penelitian}
Manfaat yang diharapkan dari penelitian ini adalah sebagai berikut:
\begin{enumerate}
    \item Memberikan pemahaman teoretis yang lebih mendalam mengenai pengisian konkurensi sebagai ukuran keterikatan multipartit sejati pada sistem tiga qubit.
    \item Menjadi referensi akademik dalam kajian keterikatan multipartit sejati, khususnya terkait klasifikasi keadaan GHZ dan W.
    \item Memberikan kontribusi konseptual bagi pengembangan studi lanjutan di bidang fisika kuantum teoretis dan informasi kuantum.
\end{enumerate}

\section{Sistematika Penelitian}
Sistematika penulisan tesis ini disusun sebagai berikut:
\begin{itemize}
    \item \textbf{Bab 1 Pendahuluan} berisi latar belakang penelitian, rumusan masalah, tujuan penelitian, batasan masalah, manfaat penelitian, serta sistematika penulisan.
    \item \textbf{Bab 2 Studi Litaratur} memuat konsep dasar terkait qubit, keterbelitan, klasifikasi keterbelitan, serta teori-teori pendukung yang relevan dengan penelitian.
    \item \textbf{Bab 3 Metodologi} menjelaskan pendekatan teoretis/metodologi penelitian yang digunakan dan langkah/diagram alir yang diterapkan dalam penelitian.
    \item \textbf{Bab 4 Hasil dan Pembahasan} menyajikan hasil analisis pengisian konkurensi pada sistem tiga qubit serta pembahasan terhadap klasifikasi keadaan GHZ dan W pada keadaan murni dan campuran.
    \item \textbf{Bab 5 Kesimpulan dan Saran} berisi kesimpulan yang diperoleh dari penelitian serta saran untuk pengembangan penelitian selanjutnya.
\end{itemize}


%\newpage
%\begin{center}
%	\topskip0pt
%	\vspace*{\fill}
%	\textit{\textbf{"Halaman ini sengaja dikosongkan"}}
%	\vspace*{\fill}
%\end{center}